\medskip

\textbf{Partie A}
v
 \begin{enumerate}
    \item $ u_1 = \dfrac{1}{2}u_0 + 10= \dfrac{1}{2} \times 30 + 10 = 15 + 10 = 25$ et $ u_2 = \dfrac{1}{2}u_1 + 10=  \dfrac{1}{2} \times 25 + 10 = 12,5 + 10 = 22,5.$

    \item Pour tout \( n \in \mathbb{N} \), $v_{n+1} = u_{n+1} - 20 = \left( \dfrac{1}{2}u_n + 10 \right) - 20 = \dfrac{1}{2}u_n - 10 = \dfrac{1}{2} \left( u_n - 20\right) = \dfrac{1}{2}v_n$.


donc la suite \( (v_n) \) est géométrique de raison \( q=\dfrac{1}{2} \) et de premier terme $v_0=u_0-20=30-20=10$.

    \item Pour tout entier $n$, $v_n=v_0 \times q^n =   10 \times \left( \dfrac{1}{2} \right)^n$

    \item Pour tout entier $n$,  $u_n = v_n + 20 = 10 \times \left( \dfrac{1}{2} \right)^n + 20 = 20 + 10 \times \left( \dfrac{1}{2} \right)^n$

    \item $\lim\limits_{n \to +\infty} \left( \dfrac{1}{2} \right)^n=0$ ( $q^n$ , avec $0<q<1$ ) donc par produit $\lim\limits_{n \to +\infty}  10 \times\left( \dfrac{1}{2} \right)^n=0$ et donc par somme  $\lim\limits_{n \to +\infty} u_n=20$ 
 
\end{enumerate}

\textbf{Partie B}
 
\begin{enumerate}
    \item $w_1= \dfrac{1}{2} \times u_0 + \dfrac{1}{2} \times 30 + 7 = \dfrac{1}{2} \times 45 + 15 + 7 = 22,5 +22 = 44,5$

    \item La fonction Python donnée ne calcule pas correctement \( w_1 \) car elle calcule $u_1$ avant de calculer $w_1$ il faut modifier l’ordre des opérations :
 
\begin{verbatim}
def suite(n):
    U = 30
    W = 45
    for i in range(1, n+1):
        W = W/2 + U/2 + 7
        U = U/2 + 10
    return W
\end{verbatim}

    \item (a) Montrons par récurrence  que pour tout entier naturel $n$ , 
   $ w_n = 10n \left( \dfrac{1}{2} \right)^n + 11 \left( \dfrac{1}{2} \right)^n + 34$

    \textbf{Initialisation :} Pour \( n = 0 \),
    $ w_0 =45 $et $  10 \times 0 \times \left( \dfrac{1}{2} \right)^0 + 11 \times \left( \dfrac{1}{2} \right)^0 + 34 =11+34= 45$

donc l'égalité est initialisée pour $n=0$.

 \textbf{Hérédité :} On suppose qu'il existe un entier $k$ tel que $w_k = 10k \left( \dfrac{1}{2} \right)^k + 11 \left( \dfrac{1}{2} \right)^k + 34$ , montrons que  $w_{k+1} = 10(k+1) \left( \dfrac{1}{2} \right)^{k+1} + 11 \left( \dfrac{1}{2} \right)^{k+1} + 34$

$ w_{k+1} = \dfrac{1}{2}w_k + \dfrac{1}{2}u_k + 7 = \dfrac{1}{2} \left[ 10k \left( \dfrac{1}{2} \right)^k + 11 \left( \dfrac{1}{2} \right)^k + 34 \right] + \dfrac{1}{2} \left[ 20 + 10 \left( \dfrac{1}{2} \right)^k \right] + 7$

donc $ w_{k+1} = 10k \left( \dfrac{1}{2} \right)^{k+1} + 11 \left( \dfrac{1}{2} \right)^{k+1} + 17 +10 +   \left( \dfrac{1}{2} \right)^{k+1}+ 7 = 10(k+1) \left( \dfrac{1}{2} \right)^{k+1} + 11 \left( \dfrac{1}{2} \right)^{k+1} + 34$

donc l'égalité est héréditaire.

\textbf{Conclusion}: l'égalité est initialisée pour $n=0$ et elle est héréditaire donc d'après le principe de récurrence, elle est vraie pour tout entier $n$.

pour tout entier naturel $n$ , $ w_n = 10n \left( \dfrac{1}{2} \right)^n + 11 \left( \dfrac{1}{2} \right)^n + 34$
 (b)  $\lim\limits_{n \to +\infty} \dfrac{10}{n}=0$ et pour tout entier $n \geqslant 4$ , $0 \leqslant 10n \left( \dfrac{1}{2} \right)^n \leqslant \frac{10}{n}$ donc d'après le théorème d'encadrement $\lim\limits_{n \to +\infty}10n \left( \dfrac{1}{2} \right)^n=0$ 

de plus $\lim\limits_{n \to +\infty}  \left( \dfrac{1}{2} \right)^n =0$ ( $q^n$ avec $0<q<1$ donc par produit $\lim\limits_{n \to +\infty} 11 \times  \left( \dfrac{1}{2} \right)^n =0$

donc par somme $\lim\limits_{n \to +\infty} w_n=34$
\end{enumerate}

\bigskip
